\newpage
\phantomsection
\label{prf:embedding-chain-compatibility}
\LRAProofFor{thm:embedding-chain-compatibility}

\begin{remark*}[Return]
\hyperref[thm:embedding-chain-compatibility]{Return to Theorem}
\end{remark*}

\begin{theorem*}[Compatibility of the Embedding Chain]
The composite $\Phi=\iota_{4}\circ\iota_{3}\circ\iota_{2}\circ\iota_{1}$ preserves $0$, $1$, addition, multiplication, and order, and sends each $n\in\mathbb{N}$ to the real number identified with $n$.
\end{theorem*}

\textbf{Professional Standard Proof.}~\\
\begin{proof}
TODO
\end{proof}

\textbf{Detailed Learning Proof.}~\\
TODO: Expand the professional proof into a detailed learning proof.

\begin{remark*}[Proof structure]
Composition of structure-preserving order embeddings is again a structure-preserving order embedding; apply each step's theorem and chase the constants and operations through the composite. The agreement claim is the statement that the diagram of canonical maps commutes.
\end{remark*}

\begin{dependencies}
\begin{itemize}
  \item \hyperref[thm:n-embeds-in-w]{$\mathbb{N}$ Embeds in $\mathbb{W}$}
  \item \hyperref[thm:q-embeds-in-r]{$\mathbb{Q}$ Embeds in $\mathbb{R}$}
\end{itemize}
\end{dependencies}
